% Vyžaduje balíčky: graphicx, float. Obrázky: analyza/analyticka_cast/obrazky/
% (\graphicspath{{analyza/analyticka_cast/obrazky/}}).
\section{Stáří zástavby a vývoj v čase}

Rok dokončení je vyplněn u 63,2\,\% aktuálních objektů. U starší zástavby
je ale třeba počítat s tím, že nejde o přesné datum: 64,4\,\% všech dat
dokončení připadá na 31.\,prosince a četnosti se výrazně hromadí na hranicích
let 1919, 1945, 1960, 1970, 1980, 1990 a 2000. Tyto hodnoty odpovídají
zařazení objektu do období, ve kterém byl postaven, a ne skutečnému roku
výstavby. Věrohodné roky jsou až od roku 2001. Starší zástavba se proto
hodnotí po obdobích a tempo výstavby po jednotlivých letech od roku 2001.

\subsection{Stáří zástavby}

Podle období dokončení (obrázek~\ref{fig:stari}) pochází 27\,\% objektů
s vyplněným datem z doby do roku 1945, 32\,\% do roku 1960 a 34\,\% z roku
2001 a později. Nejpočetnější je období 1920--1945 (12\,855 objektů),
následuje období od roku 2011 (12\,028) a 2001--2010 (10\,014). Období mají
různou délku, takže počty nejsou mezi sebou přímo srovnatelné.

\begin{figure}[H]
\centering
\includegraphics[width=0.9\textwidth]{18_stari_zastavby}
\caption{Počet objektů podle období dokončení}
\label{fig:stari}
\end{figure}

\subsection{Tempo výstavby}

Počet dokončených objektů se v letech 2001--2024 pohyboval mezi 648 (rok
2012) a 1\,429 (rok 2007) ročně (obrázek~\ref{fig:tempo}). Nejvíce objektů
bylo dokončeno v letech 2004--2007, v letech 2001--2010 v průměru asi
1\,000 ročně, od roku 2012 asi 795 ročně. Rodinné domy tvoří 64\,\% z nich,
bytové domy a garáže po 7\,\%. Roky 2025 a 2026 jsou neúplné, protože se
dokončené stavby do registru zapisují se zpožděním. Ve sledované historii
RÚIAN se za rok objeví v průměru 808 nových objektů, což odpovídá počtu
dokončených staveb zjištěnému z dat dokončení.

Graf zahrnuje pouze objekty, které v registru zůstaly. Stavby, které
mezitím zanikly, v něm nejsou, takže starší roky mohou být o něco
podhodnocené.

\begin{figure}[H]
\centering
\includegraphics[width=\textwidth]{19_tempo_vystavby}
\caption{Počet dokončených objektů podle roku dokončení a způsobu využití}
\label{fig:tempo}
\end{figure}

\subsection{Konstrukce podle stáří}

V každém období převažuje zděná konstrukce (76--93\,\%,
obrázek~\ref{fig:konstrukce}). Zřetelnou stopu zanechala panelová výstavba:
stěnové panely mají 10\,\% objektů z let 1961--1970, 17\,\% z let 1971--1980
a 16\,\% z let 1981--1990, v ostatních obdobích nejvýše 3\,\%. Celkem 79\,\%
panelových objektů pochází z let 1961--1990. Dřevěné konstrukce jsou častější
až v nejnovějším období.

\begin{figure}[H]
\centering
\includegraphics[width=\textwidth]{20_konstrukce_podle_obdobi}
\caption{Druh konstrukce podle období dokončení}
\label{fig:konstrukce}
\end{figure}

\subsection{Životnost objektů}

Historie obsahuje 22\,347 objektů, jejichž platnost skončila zánikem. Z těchto
dat ale nelze spolehlivě odvodit životnost staveb (obrázek~\ref{fig:zanik}).
O zaniklých objektech víme jen málo: využití má vyplněné 32\,\%, zastavěnou
plochu 3\,\%, podlažnost 6\,\% a datum dokončení 3\,\% (765 objektů).
Většina (71\,\%) jsou budovy s číslem evidenčním. Jde tedy převážně
o neúplné záznamy odstraněné z registru, ne o skutečně zbourané budovy. Stáří
v době zániku (medián 15 let) navíc zkreslují zástupná data typu
31.\,12.\,2000. Zánik objektu proto v aplikaci prezentujeme jako událost
v evidenci, ne jako konec fyzické existence stavby.

\begin{figure}[H]
\centering
\includegraphics[width=0.9\textwidth]{21_zanikle_objekty_vyplnenost}
\caption{Úplnost údajů u zaniklých a aktuálních objektů}
\label{fig:zanik}
\end{figure}
